% Generated draft: requires mathematical and source review.
\documentclass[11pt]{article}
\usepackage[margin=25mm]{geometry}
\usepackage{fontspec}
\setmainfont{Latin Modern Roman}
\usepackage{amsmath,amssymb,mathtools,unicode-math}
\setmathfont{Latin Modern Math}
\usepackage{enumitem,needspace,xurl,hyperref}
\setlength{\parindent}{0pt}
\setlength{\parskip}{5pt}
\setlength{\emergencystretch}{3em}
\newcommand{\N}{\mathbb{N}}
\newcommand{\Z}{\mathbb{Z}}
\newcommand{\Q}{\mathbb{Q}}
\newcommand{\R}{\mathbb{R}}
\newcommand{\C}{\mathbb{C}}
\newcommand{\F}{\mathbb{F}}
\newcommand{\A}{\mathbb{A}}
\newcommand{\PP}{\mathbb{P}}
\newcommand{\E}{\mathbb{E}}
\newcommand{\eps}{\varepsilon}
\newcommand{\dd}{\,\mathrm{d}}
\begin{document}
\section*{UnsolvedMath ID: 30003522}
\textbf{Source problem code:} catalogue entry
\subsection{Definitions and mathematical statement}
Let \(G\) be a connected semisimple Lie group with finite center, and equip \(\R\) with its usual topology and the trivial \(G\)-action. For every integer \(n\geq 0\), define the space of continuous homogeneous cochains by
\[
C^n_{\mathrm c}(G;\R)=\{f:G^{n+1}\to\R\mid f\text{ is continuous and }f(hg_0,\ldots,hg_n)=f(g_0,\ldots,g_n)\text{ for all }h,g_0,\ldots,g_n\in G\}.
\]
Its bounded subspace is
\[
C^n_{\mathrm{cb}}(G;\R)=\left\{f\in C^n_{\mathrm c}(G;\R):\sup_{(g_0,\ldots,g_n)\in G^{n+1}}|f(g_0,\ldots,g_n)|<\infty\right\}.
\]
Both complexes have differential
\[
(\delta f)(g_0,\ldots,g_{n+1})=\sum_{i=0}^{n+1}(-1)^i f(g_0,\ldots,\widehat{g_i},\ldots,g_{n+1}),
\]
where the hat denotes omission. Their cohomology groups are continuous cohomology \(H^n_{\mathrm c}(G;\R)\) and continuous bounded cohomology \(H^n_{\mathrm{cb}}(G;\R)\), respectively. These are ordinary cohomology groups, without taking closures of coboundary spaces. Inclusion of the bounded subcomplex induces the comparison map
\[
c_G^n:H^n_{\mathrm{cb}}(G;\R)\longrightarrow H^n_{\mathrm c}(G;\R).
\]
\textbf{Dupont--Monod comparison-map conjecture, as supplied.} For every connected semisimple Lie group \(G\) with finite center and every integer \(n\geq 0\), is \(c_G^n\) an isomorphism of real vector spaces?

The supplied record imposes no further restriction on compact factors or on the degree. Its mathematical formulation is consistent across the supplied statements. The report-year discrepancy and the dated literature claims described below remain unverified source qualifications.
\subsection{Short English statement}
The conjecture asks whether, for every connected semisimple Lie group with finite center, continuous bounded cohomology with trivial real coefficients agrees with continuous cohomology through the natural comparison map in every nonnegative degree. The supplied record labels the problem partially solved; this draft does not independently verify that assessment.
\subsection{Sources}
\begin{itemize}
\item Supplied local UnsolvedMath registry record, permanent ID \(30003522\), source problem code \texttt{catalogue entry}, entitled \emph{The Dupont--Monod Comparison-Map Conjecture}, in the Oberwolfach Reports Open Problems collection. The original statement attributes the conjecture by the name Dupont--Monod; no separate proposer is supplied. The proposed year is recorded as 2017.
\item Supplied report citation: \emph{Geometric Structures in Group Theory (2018)}, Oberwolfach Reports, pp.~1869--1915, DOI \texttt{10.4171/owr/2017/30}: \url{https://doi.org/10.4171/owr/2017/30}. The same reference is listed among the supplied literature sources as \emph{Oberwolfach Report 30/2017}. The citation's 2018 label and its 2017 report designation are retained without reconciliation. The background describes Oberwolfach Reports as documenting mathematical workshops and their open-problem discussions.
\item Supplied literature source: Ott, \emph{Transgression in bounded cohomology}, \url{https://arxiv.org/1811.07558}.
\item Supplied literature source: De la Cruz Mengual, \emph{Degree-three bounded cohomology}, \url{https://arxiv.org/abs/2304.00607}.
\item Literature assessment dated exactly \texttt{2026-08-21}: the supplied curation labels the problem partially solved and states that the relevant bounded-cohomology conjecture is established in low degrees and important special cases, including degree three for complex classical groups, while no general all-degree theorem is known. This is supplied dated curation derived from the OWR material, not an independently verified literature assessment; the record expressly requests independent verification of time-sensitive status and cited claims.
\item The registry labels the statement \texttt{corrected\_verified} and says its mathematical content was preserved while notation, typography, or omitted definitions were repaired for self-contained publication. This is a supplied editorial qualification, not mathematical verification performed for this draft.
\end{itemize}
\end{document}
